
\chapter*{Exposé XIII\\Regular Elements of Algebraic Groups and Lie Algebras}
\addcontentsline{toc}{chapter}{Exposé XIII. Regular Elements of Algebraic Groups and Lie Algebras}
\SGARefTarget{sga3:xiii:expose}

\begin{center}
\textit{by A. Grothendieck}
\end{center}

\section{An auxiliary lemma on varieties with operators}
\SGARefTarget{sga3:xiii:section:1}

Let \(S\) be a prescheme, let \(G\) be an \(S\)-group prescheme acting
on the left on an \(S\)-prescheme \(V\), let \(W\) be a closed
\(S\)-subprescheme of \(V\), and let \(N\) be its stabilizer in \(G\).
Thus \(N\) is the subgroup of \(G\) whose points with values in an
\(S'\) over \(S\) are the elements \(g\in G(S')\) such that
\[
  g\cdot W_{S'}=W_{S'}.
\]
We endow \((\mathrm{Sch}/S)\) with the faithfully flat
quasi-compact topology and identify \(G\), \(V\), and \(W\) with the
corresponding sheaves.\footnote{The French source prints \(M\) in this
list; the surrounding sentence and all subsequent uses require \(W\).}
We shall therefore reason in the category of sheaves on
\((\mathrm{Sch}/S)\); throughout \SGARefLink{sga3:xiii:section:1}{this section}, ``locally'' refers to
the topology just specified on \((\mathrm{Sch}/S)\).

Since \(N\) is a sheaf, consider the quotient sheaf \(G/N\). It is
immediately seen to be isomorphic to the following functor: to every
\(S'\) over \(S\), it assigns the set of subsheaves \(W'\) of
\(V_{S'}\) that are locally conjugate to \(W_{S'}\) under \(G\).
Let \(X\) be the subsheaf of \(G/N\times_S V\) whose value on an
\(S'\) over \(S\) is the set of pairs \((W',v)\), where \(W'\) is as
above and \(v\) is a section of \(W'\) over \(S'\) (hence also a
section of \(V_{S'}\) over \(S'\)). Let \(Z\) be the inverse image of
\(X\) in \(G\times_S V\). We then have \SGARefLink{sga3:xiii:equation:x}{the cartesian diagram}
\SGARefTarget{sga3:xiii:diagram:x}
\[
\tag{x}\label{sga3:xiii:equation:x}
\begin{tikzcd}[column sep=large,row sep=large]
Z \arrow[r] \arrow[d,hook] & X \arrow[d,hook,"i"] \\
G\times_S V \arrow[r] & G/N\times_S V ,
\end{tikzcd}
\]
where \(i\) is the canonical immersion, while the lower horizontal
arrow is induced by the canonical morphism \(G\to G/N\).

This morphism sends \(g\in G(S')\) to the subsheaf
\(g\cdot W_{S'}\) of \(V_{S'}\). Consequently, \(Z(S')\) is the set
of pairs \((g,v)\in G(S')\times V(S')\) such that
\(v\in (g\cdot W_{S'})(S')\). It follows that \(Z\) is isomorphic to
the sheaf \(G\times_S W\), via the isomorphism
\[
  G\times_S W \xrightarrow{\sim} Z,\qquad
  (g,w)\longmapsto(g,g\cdot w).
\]
Thus \SGARefLink{sga3:xiii:equation:x}{the preceding cartesian diagram} yields \SGARefLink{sga3:xiii:equation:xx}{the cartesian diagram}
\SGARefTarget{sga3:xiii:diagram:xx}
\[
\tag{xx}\label{sga3:xiii:equation:xx}
\begin{tikzcd}[column sep=large,row sep=large]
G\times_S W \arrow[r,"q"] \arrow[d,"\lambda"'] &
  X \arrow[d,"i"] \\
G\times_S V \arrow[r] & G/N\times_S V
\end{tikzcd}
\]
where
\[
  \lambda(g,w)=(g,g\cdot w),\qquad
  q(g,w)=(\bar g,g\cdot w),
\]
where \(\bar g\) denotes the image of \(g\) under the canonical map
\(G(S')\to(G/N)(S')\).

Diagram \(\SGARefLink{sga3:xiii:equation:x}{\mathrm{(x)}}\) also shows that \(Z\to X\) makes \(Z\) a
principal bundle over \(X\), with \(N\) acting on the right by
\[
  (g,v)\cdot n=(gn,v).
\]
Accordingly, in diagram \(\SGARefLink{sga3:xiii:equation:xx}{\mathrm{(xx)}}\), the morphism
\(q:G\times_S W\to X\) makes \(G\times_S W\) a principal bundle over
\(X\), with right \(N\)-action
\[
  (g,w)\cdot n=(gn,n^{-1}\cdot w).
\]

The principal morphisms just introduced are summarized by
\SGARefTarget{sga3:xiii:diagram:d}
\[
\tag{D}\label{sga3:xiii:equation:d}
\begin{tikzcd}[column sep=large,row sep=large]
& G\times_S W \arrow[d,"q"] \arrow[dr,"\varphi"] & \\
G/N &
X \arrow[l,"p"] \arrow[r,"\psi"] \arrow[d,hook,"i"] &
V \\
& G/N\times_S V
  \arrow[ul,dashed,"\operatorname{pr}_1"]
  \arrow[ur,dashed,"\operatorname{pr}_2"'] &
\end{tikzcd}
\]
where
\[
  \psi=\operatorname{pr}_2\circ i,\qquad
  \varphi=\psi\circ q,\qquad
  \varphi(g,w)=g\cdot w.
\]

If \(v\) is a section of \(V\) over \(S\), let \(X_v\) be its inverse
image in \(X\) under \(\psi\). For every \(S'\) over \(S\),
\(X_v(S')\) is the set of subsheaves \(W'\) of \(V_{S'}\) that are
locally conjugate to \(W_{S'}\) under \(G\) and contain the section
\(v_{S'}\) of \(V_{S'}\). The inverse image of \(v\) in
\(G\times_S W\) under \(\varphi\) is isomorphic to the subsheaf
\(M_v\) of \(G\) defined by
\[
  M_v(S')=
  \bigl\{g\in G(S')\mid v_{S'}\in(g\cdot W_{S'})(S')\bigr\}
  =
  \bigl\{g\in G(S')\mid g^{-1}v_{S'}\in W(S')\bigr\}.
\]
If \(v\) is a section of \(W\), and not merely of \(V\), then
\(M_v\) evidently contains \(N\).

None of these explicit descriptions used the representability of
\(G\), \(V\), or \(W\) (nor even the fact that the site in question is
defined in terms of preschemes). Suppose now, however, that \(N\) is
representable, faithfully flat, and quasi-compact over \(S\), and that
\(G/N\) is representable. When \(S\) is the spectrum of a field and
\(G\) is of finite type over that field, this hypothesis is necessarily
satisfied (VIII~6 and VI\textsubscript{B}.11.18). Using the cartesian
diagram \(\SGARefLink{sga3:xiii:equation:xx}{\mathrm{(xx)}}\), faithfully flat quasi-compact descent, and
the fact that \(Z\to G\times_S V\) is a closed immersion, one sees that
\(X\) is representable: it is obtained by descending the closed
subprescheme \(Z\) of \(G\times_S V\) along the faithfully flat
quasi-compact morphism
\[
  G\times_S V\longrightarrow G/N\times_S V.
\]
Thus diagram \(\SGARefLink{sga3:xiii:equation:d}{\mathrm{(D)}}\) is a diagram of morphisms of
preschemes over \(S\).
